\documentclass[12pt]{article}

% ------------------------------------------------------------
% Encoding and Fonts
% ------------------------------------------------------------
\usepackage[T1]{fontenc}
\usepackage[utf8]{inputenc}
\usepackage{lmodern}

% ------------------------------------------------------------
% Mathematics Packages
% ------------------------------------------------------------
\usepackage{amsmath, amssymb, amsthm, mathtools}
\usepackage{bm}
\usepackage{bbm}
\usepackage{physics}
\usepackage{enumitem}
\usepackage{geometry}
\usepackage{xurl}
\usepackage{csquotes}
\usepackage{hyperref}

% ------------------------------------------------------------
% Layout
% ------------------------------------------------------------
\geometry{margin=1in}
\setlength{\parindent}{0pt}
\setlength{\parskip}{6pt}

% ------------------------------------------------------------
% Bibliography
% ------------------------------------------------------------
\usepackage[
    backend=biber,
    style=authoryear,
    natbib=true,
    maxbibnames=10,
    maxcitenames=2,
    doi=false,
    url=true
]{biblatex}

\addbibresource{references.bib}

% ------------------------------------------------------------
% Theorem Environments
% ------------------------------------------------------------
\theoremstyle{plain}
\newtheorem{theorem}{Theorem}[section]
\newtheorem{lemma}[theorem]{Lemma}
\newtheorem{proposition}[theorem]{Proposition}
\newtheorem{corollary}[theorem]{Corollary}

\theoremstyle{definition}
\newtheorem{definition}[theorem]{Definition}
\newtheorem{assumption}[theorem]{Assumption}

\theoremstyle{remark}
\newtheorem{remark}[theorem]{Remark}

% ------------------------------------------------------------
% Custom Math Operators
% ------------------------------------------------------------
\DeclareMathOperator{\Ric}{Ric}
\DeclareMathOperator{\vol}{vol}
\DeclareMathOperator{\Span}{Span}
\DeclareMathOperator{\rank}{rank}
\DeclareMathOperator{\Hess}{Hess}
\DeclareMathOperator{\Div}{div}

% ------------------------------------------------------------
% Notation Macros
% ------------------------------------------------------------
\newcommand{\ThetaSpace}{\Theta}
\newcommand{\ModelMap}{f_\theta}
\newcommand{\ParamJac}{J_\theta}
\newcommand{\Stratum}{\mathcal{S}}
\newcommand{\Realization}{\Phi}
\newcommand{\FreeEnergy}{\mathcal{F}}
\newcommand{\LamEnergy}{\mathcal{L}_{\mathrm{lam}}}
\newcommand{\SI}{\mathsf{SI}}
\newcommand{\Null}{\mathcal{N}}
\newcommand{\TartanOp}{\mathsf{T}}
\newcommand{\LamFlow}{\mathsf{L}}

% ------------------------------------------------------------
% Title
% ------------------------------------------------------------
\title{\Large\bfseries
Recursive Agency and the Geometry of the Relativistic Scalar--Vector Plenum:\\
A Unified Field Theory of Representation, Semantics, and Computational Life
}
\author{Flyxion}
\date{\today}

\begin{document}

\maketitle
\tableofcontents
\newpage

% ------------------------------------------------------------
% ABSTRACT
% ------------------------------------------------------------
\section*{Abstract}

Intelligent systems---biological, artificial, and theoretical---inhabit 
manifolds of possibility whose structure governs what can be represented, 
predicted, or enacted. The Relativistic Scalar--Vector Plenum (RSVP) theory
proposes a unifying field framework in which representation, inference, 
semantic flow, and agency emerge from the geometry and dynamics of coupled 
scalar, vector, and entropy fields. Deep learning architectures, cognitive 
agents, variational flows, and physical systems reveal a shared geometric 
logic: they evolve along stratified manifolds, regularize via lamphrodynamic 
smoothing, and preserve stability by maintaining coherence across scales. 

This monograph synthesizes the mathematics of representation manifolds, 
stratified geometry, fiber symmetry, semantic fields, multiscale inference, 
free-energy flows, lamphrodynamic complexity, predictive silence, and the 
constraints of mortal computation. We develop a unified geometric language 
for recursive agency: systems that continually reshape the very manifolds 
in which they operate. The result is a coherent mathematical cosmology of 
intelligence, grounded in Riemannian geometry, operator theory, category 
theory, renormalization, thermodynamics, and dynamical systems.

\newpage

% ------------------------------------------------------------
% PART I: REPRESENTATION GEOMETRY
% ------------------------------------------------------------
\part{Representation Geometry}

\section{Representation Manifolds}

Modern learning systems inhabit parameter spaces $\ThetaSpace$ of enormous 
dimension, yet their functional degrees of freedom lie on low-dimensional, 
curvature-rich manifolds induced by realization maps. These 
\emph{representation manifolds} encode the expressive potential of a model 
independently of any particular parameterization.

Let $\ModelMap(\cdot;\theta)$ be a parametric model. The representation 
manifold is the image:
\[
M = \Realization(\ThetaSpace) \subset \mathcal{H},
\]
or in discrete form $M \subset \mathbb{R}^N$. Although $\ThetaSpace$ may 
be linear, $M$ is almost always curved, stratified, and rife with symmetry:
the true geometry of representation is nonlinear even when the parameters are not.

The Jacobian
\[
\ParamJac_\theta : T_\theta \ThetaSpace \to T_{\Realization(\theta)} M
\]
determines the local representational freedom at $\theta$. Its rank varies 
across strata of $\ThetaSpace$, producing a Whitney-stratified structure 
reflecting different modes of expressivity.

\section{Rank Stratification and Capacity}

Define the rank strata:
\[
\Stratum_r = \{\theta \in \ThetaSpace : \rank \ParamJac_\theta = r\}.
\]
Each stratum corresponds to a distinct regime of representational capacity.

In deep networks, strata typically correspond to patterns of activated or 
inactive ReLUs, degenerate convolutional kernels, collapsed attention heads, 
or linear dependencies among internal features. These structures are 
semi-algebraic, and thus admit a finite Whitney stratification 
\citep{goresky1988stratified, mather1973stratifications}.

The geometry is not static: as learning proceeds, updates traverse strata, 
create or annihilate symmetries, and restructure expressivity. Capacity is 
therefore not a scalar resource but a geometric field distributed across $\ThetaSpace$.

\section{Fiber Geometry and Symmetry}

Even within a rank-constant region, many parameters correspond to the same 
representational function. The fiber:
\[
F(\theta) = \Realization^{-1}(\Realization(\theta))
\]
encodes these redundancies. Its dimension is:
\[
\dim F(\theta) = \dim \ThetaSpace - \rank \ParamJac_\theta.
\]

Symmetry groups $G$ acting on parameters without changing representations 
embed their orbits inside fibers. These include permutation symmetries, 
scaling groups in ReLU networks, and continuous gauge freedoms. Such symmetries 
reflect flat directions in the loss landscape and are geometrically essential 
to understanding generalization and agency.

\section{Vertical and Horizontal Dynamics}

Within a rank-constant stratum, tangent vectors split into:

- vertical directions $V_\theta = \ker \ParamJac_\theta$ preserving the model’s output,  
- horizontal directions modifying output locally.

This induces a partial connection on $\ThetaSpace$, enabling analysis of 
learning rules in geometric terms. Gradient-based learning generally mixes 
horizontal motion (improving performance) with vertical drift (exploration 
of flat directions). Agency emerges when the system modulates this decomposition.

\newpage

% ------------------------------------------------------------
% PART II: SEMANTIC FIELDS AND LAMPHRODYNAMICS
% ------------------------------------------------------------
\part{Semantic Fields and Lamphrodynamic Dynamics}

\section{Semantic Fields on Representation Manifolds}

A representation manifold $M$ is not static: it carries semantic fields that
encode predictions, values, features, or internal potentials. A semantic field
is a map:
\[
f : M \to \mathbb{R}^k,
\]
whose evolution reflects the internal dynamics of learning or inference.
We interpret semantic fields as generalized coordinates of meaning. Their 
dynamics formalize how an intelligent system stabilizes, refines, or reshapes 
its internal semantics.

Given a functional
\[
\mathcal{E}(f) 
= \int_M \left( V(f(x)) + \frac{1}{2} \abs{\nabla f(x)}^2 \right) d\vol_g(x),
\]
the induced gradient flow is:
\[
\partial_t f_t 
= - \frac{\delta \mathcal{E}}{\delta f}(f_t)
= \Delta f_t - \nabla V(f_t).
\]
This is a reaction–diffusion equation on $M$, whose qualitative behavior depends 
crucially on the curvature of $M$, the shape of $V$, and the stratified structure
inherited from the realization map.

\subsection{Semantic Coherence and Energetic Descent}

The functional $\mathcal{E}$ plays the role of a semantic free energy: it 
quantifies excess tension or irregularity in the system’s predictions or 
internal potentials. Along the gradient flow,
\[
\frac{d}{dt} \mathcal{E}(f_t)
= - \int_M \abs{\partial_t f_t}^2\, d\vol_g \le 0.
\]
Semantic flows thus dissipate energy and increase coherence, smoothing internal
representations until only meaningful structure remains.

The rate and pattern of semantic smoothing depend on the geometric properties 
of $M$. High curvature accelerates dissipation in some regions and retards it 
in others; singular strata impede or redirect flows; symmetry or fiber structure
induces degeneracy or invariance. Learning is geometrically entangled with 
the manifold on which meaning evolves.

\section{Lamphrodynamic Energy and Higher-Order Regularization}

Classical diffusive smoothing, such as the heat flow, suppresses local variation.
But intelligent systems require a balance: they must smooth noise while preserving 
or amplifying semantically important structure. This motivates lamphrodynamic
regularization.

Define the lamphrodynamic energy:
\[
\LamEnergy(f)
= \int_{\mathbb{R}^d} \left( 
      \alpha \abs{\nabla_x f(x)}^2
    + \beta \abs{\Hess_x f(x)}^2
  \right) d\mu(x),
\]
where $\mu$ is an input distribution and $\alpha, \beta \ge 0$.

The $H^2$ term $\abs{\Hess f}^2$ penalizes high curvature in the mapping 
$f : \mathbb{R}^d \to \mathbb{R}^k$, inducing a bi-harmonic smoothing flow:
\[
\partial_t f
= 2\alpha \Delta f - 2\beta \Delta^2 f.
\]
This fourth-order PDE is more stable than heat flow: it flattens jaggedness
but preserves low-curvature semantic modes. In RSVP dynamics, this corresponds
to lamphrodynamic smoothing: the dissipation of semantic turbulence.

\section{Existence and Uniqueness of Lamphrodynamic Flows}

Higher-order parabolic equations admit global smooth solutions under
standard boundary and regularity conditions. Following the theory in 
\citep{evans2010partial, davies1980oneparameter}, the flow
\[
\partial_t f = - \frac{\delta \LamEnergy}{\delta f}
\]
is well-posed in $H^2$ with solutions:
\[
f \in C([0,\infty); H^2) \cap L^2_{\mathrm{loc}}([0,\infty); H^4).
\]
Moreover, lamphrodynamic energy decays monotonically:
\[
\frac{d}{dt} \LamEnergy(f_t)
= - \int \abs{\partial_t f_t}^2\, d\mu \le 0.
\]
Lamphrodynamics thereby enforces long-term stability of semantic fields even 
on curved or stratified representation manifolds.

\section{Predictive Conductance and Curvature}

Lamphrodynamic flows interact deeply with curvature. On positively curved
manifolds, diffusion is amplified; on negatively curved manifolds, it disperses.
The bi-harmonic term partially neutralizes curvature effects, enabling coherent
semantic propagation across regions of varying geometric sharpness.

In RSVP theory, lamphrodynamics mediates how semantic tension distributes across
scalar and vector fields, smoothing excessive divergence, dissipating localized
semantic vortices, and producing stable attractor structures aligned with global
geometric constraints.

\newpage

% ------------------------------------------------------------
% PART III: MULTISCALE TRAJECTORIES AND TARTAN GEOMETRY
% ------------------------------------------------------------
\part{Multiscale Trajectories and TARTAN Geometry}

\section{Trajectory Fields and Multiscale Memory}

Modern inference systems operate not only on single states but on trajectories.
A sequence of hidden states $\{z_t\}_{t=0}^T \subset \mathbb{R}^d$ forms a 
\emph{trajectory field}. To stabilize these trajectories, we introduce 
multiscale aura fields:
\[
a_t 
= \frac{\sum_{k=0}^{K-1} \exp\!\left( -\frac{\abs{z_t - b_k}^2}{2\sigma^2} \right) b_k}
       {\sum_{k=0}^{K-1} \exp\!\left( -\frac{\abs{z_t - b_k}^2}{2\sigma^2} \right)}.
\]
Here $\{b_k\}$ is a buffer of prior states representing the ``memory'' or 
``history'' of the system. The smoothed update:
\[
\tilde{z}_t 
= z_t - \lambda (z_t - a_t)
\]
implements a contraction toward semantically relevant past states.

\subsection{Nonexpansiveness and Stability}

For $0 < \lambda \ll 1$, the update map $z_t \mapsto \tilde{z}_t$ is 
nonexpansive:
\[
\abs{\tilde{z}_t - \tilde{z}_t'}
\le \abs{z_t - z_t'}.
\]
This ensures stability of the trajectory in the presence of noise or 
model uncertainty. TARTAN (Trajectory-Aware Recursive Tiling with Annotated
Noise) exploits this to maintain consistent multiscale structure.

\section{Aura Geometry and Recursive Smoothing}

Aura fields behave as geometric attractors: they encode locality in both 
space and semantic coherence. The Gaussian weights define a local chart of
the trajectory manifold, and contraction toward these charts yields 
multiscale regularization.

From a geometric viewpoint, the TARTAN operator $\TartanOp$ defines a 
flow on trajectory space that contracts noisy or chaotic directions 
while preserving or amplifying meaningful ones. It is a discrete, 
computational implementation of a lamphrodynamic smoothing flow 
on temporal manifolds.

\section{Semantic Turbulence and Multiscale Coherence}

Semantic turbulence arises when predictions or internal representations fluctuate 
across incompatible scales. Multiscale coherence requires that local inference 
align with long-range semantic structure.

TARTAN geometry formalizes this requirement. A consistent trajectory satisfies:
\[
D_{\mathrm{KL}}(\tilde{z}_{t+k} \Vert z_t) \text{ small for all relevant } k.
\]
This expresses global coherence: future semantic states must be reachable from 
past ones through low-energy transformations.

\section{Recursive Inference and Trajectory Attractors}

Recursive agency emerges when the system not only smooths its trajectory but 
modifies the rules by which smoothing occurs. The TARTAN operator may itself 
adapt across scales, creating a hierarchy of attractors:
\[
\TartanOp_0 \leadsto \TartanOp_1 \leadsto \cdots \leadsto \TartanOp_L.
\]
Each scale encodes a distinct level of abstraction or semantic granularity.
This multilevel structure parallels renormalization in physics and 
multiresolution analysis in harmonic analysis.

\newpage

% ------------------------------------------------------------
% PART IV: STRATIFIED FREE ENERGY AND PREDICTIVE SILENCE
% ------------------------------------------------------------
\part{Stratified Free Energy and Predictive Silence}

\section{Variational Geometry on Stratified Manifolds}

Inference in generative models is governed by the variational free energy
functional:
\[
\FreeEnergy(\theta,\phi;x)
= \mathbb{E}_{q_\phi(z\mid x)}
  \left[
    \log q_\phi(z\mid x) - \log p_\theta(x,z)
  \right].
\]
Classically, optimization proceeds by descending this functional. But in 
geometrically nontrivial parameter spaces, $\FreeEnergy$ becomes a field
defined on a stratified manifold.

Let $\ThetaSpace \times \Phi$ be stratified according to the rank of the 
joint realization map:
\[
(\theta,\phi) \mapsto (p_\theta, q_\phi).
\]
Within each stratum, $\FreeEnergy$ is smooth; across strata, its gradient
flow interacts with changing ranks, emerging symmetries, and degeneracies.
This induces a \emph{stratified free-energy geometry}.

\subsection{Critical Points in Stratified Spaces}

A point $(\theta,\phi)$ is a critical point if the restriction of $\FreeEnergy$
to the stratum containing it has vanishing gradient:
\[
\nabla \FreeEnergy
\big\vert_{\Stratum_\alpha} = 0.
\]
Under standard conditions \citep{goresky1988stratified}, each stratum 
admits Morse-theoretic analysis. Critical points correspond to stable or 
unstable modes of inference, encoded by the curvature of $\FreeEnergy$ 
along directions tangent to the stratum.

\subsection{Stratified Descent Dynamics}

Gradient descent follows:
\[
\dot{\theta}_t = -\nabla_\theta \FreeEnergy,
\qquad
\dot{\phi}_t = -\nabla_\phi \FreeEnergy.
\]
As trajectories cross between strata:

- Rank may increase or decrease.  
- Symmetries may appear or vanish.  
- Fiber dimensions may expand or collapse.  
- Degenerate directions may unfold or fold shut.  

This induces geometric ``phase transitions'' in inference---turning points
where the system reorganizes its representational and inferential modes.

\section{Predictive Silence as Harmonicity}

Predictive silence refers to a steady state in which the system emits no new 
semantic corrections except where externally forced. Formally, a semantic map 
$u : M \to \mathbb{R}^k$ achieves predictive silence when:
\[
\Delta u = 0.
\]
Thus $u$ is harmonic. 

\subsection{Boundary-Driven Semantics}

Given fixed boundary conditions $u \vert_{\partial M}$, classical elliptic 
theory \citep{evans2010partial} guarantees:

\begin{theorem}[Existence and Uniqueness of Harmonic Semantics]
Let $M$ be a compact Riemannian manifold with smooth boundary, and let 
$f : \partial M \to \mathbb{R}^k$ be smooth. Then there exists a unique 
harmonic map $u : M \to \mathbb{R}^k$ such that $u\vert_{\partial M} = f$.
\end{theorem}

Thus predictive silence corresponds to a global semantic equilibrium shaped 
entirely by constraints at the informational boundary.

\subsection{Harmonic Agency}

Harmonic maps minimize the Dirichlet energy:
\[
E(u) = \frac{1}{2} \int_M \abs{\nabla u}^2 \, d\vol_g.
\]
In agency-theoretic terms, harmonicity represents a stable state of 
maximally efficient internal representation: all superfluous semantic 
turbulence has been smoothed away, leaving only essential structure.

Predictive silence therefore emerges as the endpoint of lamphrodynamic and 
semantic descent dynamics: a condition of informational minimality constrained
only by external data.

\section{Harmonic Limits of Learning Dynamics}

Under broad conditions (convex potentials, appropriate curvature bounds), 
solutions to the semantic flow:
\[
\partial_t f_t = \Delta f_t - \nabla V(f_t)
\]
converge to a harmonic map when $V$ is constant or negligible. Even when 
$V$ introduces nonlinearities, the long-term dynamics resemble harmonic 
projections smoothed by lamphrodynamic regularization.

This provides a geometric explanation for generalization: the system converges 
to a minimal-complexity extension of boundary-defined semantics.

\section{Predictive Precision and Entropic Curvature}

Predictive silence does not imply predictive weakness. A harmonic map can be 
sharp and precise when boundary conditions are strong. What disappears is 
unmotivated internal noise.

This is compatible with the free-energy principle \citep{friston2010free}:
minimizing variational free energy corresponds to minimizing prediction error
and model complexity---harmonicity is the geometric manifestation of this 
principle under smoothness constraints.

In RSVP, this equilibrium arises when scalar–vector fields achieve coherence:
entropy gradients align with curvature, semantic flux dissipates, and the 
system rests in a state of constrained anticipation.

\newpage

% ------------------------------------------------------------
% PART V: CONSTRAINTS, OPTIMAL CONTROL, AND COMPUTATIONAL MORTALITY
% ------------------------------------------------------------
\part{Constraints, Optimal Control, and Computational Mortality}

\section{Computational Life and Energy Budgets}

Every computational agent operates under finite resources: energy, time,
memory, and representational precision. Let $E(t)$ denote available energy,
with consumption rate $P(t)$ and replenishment $I(t)$:
\[
\frac{dE}{dt} = -P(t) + I(t).
\]
Computation becomes impossible once $E(t)$ drops below $E_{\mathrm{min}}$,
the minimal energy required to maintain system coherence.

\begin{definition}[Mortality Time]
Given a policy $\pi$, the computational mortality time is:
\[
T_{\mathrm{mort}}(\pi)
= \inf\{ t : E(t) \le E_{\mathrm{min}} \}.
\]
\end{definition}

The central challenge is to maximize performance subject to survival: 
computation is not free.

\section{Optimal Control of Computational Flow}

Let $x_t$ represent internal semantic or representational state, with control 
$u_t$ modulating computational activity. Dynamics follow:
\[
\dot{x}_t = F(x_t,u_t),
\qquad
\dot{E}_t = -P(x_t,u_t) + I_t.
\]
We define an objective:
\[
J(\pi)
= \mathbb{E}\left[
    \int_0^{T_{\mathrm{mort}}(\pi)}
       \ell(x_t,u_t)\, dt
  \right],
\]
where $\ell$ trades task performance against energy consumption.

The associated Hamilton–Jacobi–Bellman equation 
\citep{bertsekas2005dynamic, khalil2002nonlinear} governs optimal policies.

\section{Computational Drift and Renormalized Precision}

As resources diminish, fine-grained semantic distinctions become unsustainable.
Let $\Theta_\ell$ denote the representational manifold at scale $\ell$.
Available scales satisfy:
\[
\Theta_0 \supset \Theta_1 \supset \cdots \supset \Theta_L,
\]
with finer scales lost as computational entropy increases.

This induces \emph{semantic drift}: the agent's priors and representations 
collapse toward coarse-grained structures. Agency shrinks as expressive 
capacity contracts.

\section{Renormalization of Priors and Coherence Across Scales}

Let $p_\ell(\theta)$ be the prior over $\Theta_\ell$. Renormalization maps 
fine-scale priors to coarse ones:
\[
p_{\ell+1}(\theta_{\ell+1})
\propto
\int_{\pi_\ell^{-1}(\theta_{\ell+1})}
  e^{-C_\ell(\theta)}\, d\mu_\ell(\theta),
\]
where $C_\ell$ encodes energy, curvature, and entropy costs.

Renormalized priors remain Gibbsian:
\[
p_{\ell+1}(\theta_{\ell+1})
\propto
e^{-\beta_{\ell+1} C_{\ell+1}(\theta_{\ell+1})}.
\]

Coherence across scales requires:
\[
D_{\mathrm{KL}}(
  p_{\ell+1}
  \,\Vert\, 
  \mathcal{R}_\ell(p_\ell)
) < \varepsilon_\ell.
\]
Failure of this inequality constitutes \emph{agency collapse}: fine-scale 
updates conflict with coarse-scale feasibility.

\section{Semantic Horizons and the Limits of Inference}

Define the semantic horizon $\ell_{\max}$ as the finest sustainable scale:
\[
\ell_{\max}
= \sup\left\{
\ell : \beta_\ell \abs{\nabla C_\ell} \le E_{\max}
\right\}.
\]
This horizon determines the upperbound of representational precision and 
cognitive depth. Beyond it, distinctions are smoothed away by lamphrodynamic,
entropic, or energetic collapse.

Computational mortality describes the global death of agency; semantic 
horizons describe its gradual thinning.

\newpage

% ------------------------------------------------------------
% PART VI: UNIFIED FORMALISM OF THE RELATIVISTIC SCALAR--VECTOR PLENUM
% ------------------------------------------------------------
\part{Unified Formalism of the Relativistic Scalar--Vector Plenum}

\section{The Scalar--Vector Field Architecture}

The Relativistic Scalar--Vector Plenum (RSVP) proposes that intelligence,
representation, and inference are manifestations of the dynamics of three
interacting fields:

- a scalar potential field $\Phi$,  
- a vector flow field $\mathbf{v}$,  
- an entropy field $S$.

Together, these define a geometric substrate in which semantic structures
emerge, stabilize, or decay. RSVP is not a conventional physical theory nor
a metaphor: it is a geometric framework capable of expressing the deep 
structural relationships shared across learning systems, dynamical systems,
and semantic computation.

\subsection{Constitutive Equations}

At the heart of RSVP is a coupled system of PDEs representing transport,
dissipation, and recursive self-modulation:
\[
\partial_t \Phi = -\kappa \abs{\mathbf{v}}^2 \Phi
                  + \alpha \abs{\nabla\cdot \mathbf{v}}
                  + \Delta \Phi
                  - \beta (\Phi - \Phi_0),
\]
\[
\partial_t \mathbf{v} 
= -\nabla \Phi
  - \gamma\, \mathbf{v}
  + \nu \Delta \mathbf{v}
  + \eta\, \nabla S,
\]
\[
\partial_t S
= \sigma \Delta S
  + \rho \abs{\nabla \Phi}
  - \lambda S.
\]

These equations represent capacity consumption through the term $-\kappa \lvert \mathbf{v} \rvert^{2} \Phi$, which expresses the manner in which vector-field activity depletes the scalar potential and thereby regulates the availability of organizational resources. They further articulate the generation of entropy through the divergence of the vector field, a mechanism that captures the subtle interaction between informational dissipation and the geometric structure of flow. In addition, the equations incorporate the effects of vector dissipation arising from damping, viscous diffusion, and scalar-driven attraction, each of which contributes to the attenuation, redistribution, and stabilization of momentum within the medium. Finally, they formulate a semantic entropy balance that couples meaning, curvature, and flow, thereby illuminating the deep structural relationship between coherence, geometric deformation, and the dynamics of uncertainty in complex systems.


The constants $(\kappa,\alpha,\beta,\gamma,\nu,\eta,\sigma,\rho,\lambda)$ 
encode medium-specific properties.

\section{Semantic Operators and the RSVP Metric}

To express inference and representation within the RSVP substrate, define the
metric:
\[
g = \Phi^{-1} \delta_{ij} + \tau (v_i v_j),
\]
where $\tau$ is a coupling parameter. This metric induces:

- curvature $K$ depending on scalar density and vector alignment,  
- a Laplacian $\Delta_g$ governing semantic diffusion,  
- a natural geodesic structure encoding minimal semantic change.

Semantic fields evolve according to:
\[
\partial_t f = -\Delta_g f + \nabla_g S,
\]
capturing the interaction between meaning and entropy.

\section{Derived Semantics and Higher-Order Structure}

Beyond the base fields $(\Phi,\mathbf{v},S)$, RSVP admits derived semantic
objects. For a semantic field $u$, define its derived form:
\[
Du = (\nabla u, \Delta u, \Delta^2 u, \ldots ),
\]
creating a tower of higher-order invariants. Semantic structure often resides
not in $u$ itself but in these derived quantities—parallel to derived 
categories in algebraic geometry \citep{bredon1997sheaf, connes1994noncommutative}.

The lamphrodynamic energy emerges naturally in this hierarchy as the 
second-order part.

\section{Operator Algebras and Semantic Cohomology}

Let $\mathcal{A}$ be the algebra generated by:

- multiplication by scalar potentials,  
- Lie derivatives along $\mathbf{v}$,  
- Laplacian and bi-Laplacian operators,  
- entropy-gradient operators.

Then $\mathcal{A}$ forms a noncommutative algebra representing the system’s
semantic operations. Elements of $\mathcal{A}$ encode permissible semantic
transformations.

The cohomology of this algebra,
\[
H^\bullet(\mathcal{A}),
\]
encodes invariants of meaning: semantic structures preserved across all 
internal transformations. These invariants define semantic ``fixed points'':
stable conceptual entities in the agent’s internal world model.

\section{Stratified Agency and Higher Functors}

Agency is modeled as a functor:
\[
\mathcal{A} : \mathsf{Ctx} \to \mathsf{Act},
\]
mapping semantic contexts to actions. But RSVP introduces a more refined
picture: contexts and actions lie in stratified categories with varying 
degrees of representational capacity.

Define $\mathsf{Ctx}_r$ as the category of contexts defined on stratum 
$\Stratum_r$. The composite functor:
\[
\mathcal{A}_r : \mathsf{Ctx}_r \to \mathsf{Act}_r
\]
represents agency at a given representational rank.

Recursive agency arises when:
\[
\mathcal{A} : \mathsf{Ctx} \to \mathsf{Act}
\qquad\text{and}\qquad
\mathcal{A} \in \mathsf{Ctx}
\]
simultaneously: the agent acts on the category that defines its own 
action–context mappings.

\section{Renormalization, Semantic Flow, and Referential Coherence}

Semantic content evolves across scales. Let $S_\ell$ be the entropy field at 
scale $\ell$. Renormalization maps:
\[
(\Phi_\ell, \mathbf{v}_\ell, S_\ell)
\leadsto
(\Phi_{\ell+1}, \mathbf{v}_{\ell+1}, S_{\ell+1}),
\]
via coarse-graining operators consistent with the RSVP equations.

Coherence requires:
\[
D_{\mathrm{KL}}(
p_{\ell+1}(u)
\Vert 
\mathcal{R}_\ell(p_\ell)(u)
)
\le \varepsilon_\ell.
\]
This ensures that meaning does not fracture across scales.

\section{Fixed Points and Stable Semantic Attractors}

Under renormalization, certain field configurations remain invariant:
\[
(\Phi,\mathbf{v},S)
= \mathcal{R}_\ell(\Phi,\mathbf{v},S).
\]

These fixed points represent stable regimes of intelligence, characterized by coherence, energetic efficiency, and semantic consistency. In such regimes the scalar potential settles into harmonic configurations, the vector field approaches a divergence-free form that reflects a balanced and incompressible flow of influence, and the entropy field attains a stationary profile in which production and dissipation are held in equilibrium. Taken together, these properties define a state of predictive silence, in which the agent refrains from generating spurious or unstable predictions yet does not collapse into trivial inactivity. Rather, the system resides in a dynamically stable configuration that preserves meaningful structure while minimizing unnecessary informational or energetic expenditure.

\section{Unified View: Representation, Dynamics, and Conscious Computation}

RSVP achieves a synthesis of the major themes developed throughout the preceding parts by integrating the geometry of representation, the dynamics of semantic flows and lamphrodynamic regularization, the variational structure of free-energy functionals, and the energetic as well as computational constraints inherent to any finite information-processing system. Through this integration, it becomes possible to view intelligent behavior as the sustained maintenance of coherent semantic structure across multiple scales, where such structure must be continually preserved against the dissipative forces imposed by noise, curvature, and the limitations of available resources. 

At its highest level of abstraction, intelligence arises from the capacity of the system to regulate not only the evolution of its scalar, vector, and entropy fields, but also the laws that govern their interaction. This recursive modulation of both state and rule constitutes the essence of recursive agency. Within this framework, conscious computation can be understood as the geometric equilibrium achieved when meaning, prediction, and action become mutually self-consistent. Such an equilibrium reflects a state in which semantic coherence is preserved, uncertainty is minimized in a principled manner, and the dynamics of the plenum support a stable yet adaptable configuration capable of sustaining goal-directed inference and behavior.

\newpage

% ------------------------------------------------------------
% APPENDICES
% ------------------------------------------------------------
\appendix

\section{Differential-Geometric Preliminaries}
\label{app:differential-geometry}

This appendix recalls several foundational concepts from differential geometry that are used throughout the monograph. Although all results presented here are classical, they are included to clarify notation and to emphasize the geometric structure underlying recursive agency. The goal is to foreground the role of smooth manifolds, Riemannian metrics, curvature, stratification, and semi-algebraic structure in the analysis of representation and inference.

Let $(M,g)$ be a smooth, finite-dimensional Riemannian manifold. The Levi--Civita connection associated with $g$ is denoted by $\nabla$, and the Laplace--Beltrami operator by $\Delta = \mathrm{div}\circ\nabla$. The Riemann curvature tensor is written $R$, the Ricci tensor by $\Ric$, and the scalar curvature by $R_g$. The Bochner identity, already utilized in the main text, implies deep connections between harmonicity, curvature, and diffusion. Such identities justify the use of curvature bounds in establishing the well-posedness and stability of semantic flows.

In the context of RSVP, the scalar field $\Phi$ modifies the effective factor by which distances contribute to the local geometry. When $\Phi$ is large, the region becomes representationally dense; when $\Phi$ is small, the region is sparse. The induced metric $g_\Phi = \Phi^{-1} g$ amplifies this effect. Because $\Phi$ may vary across strata, the geometry of $M$ becomes stratified, and the standard comparison principles require careful adaptation. This motivates the stratified approach to curvature and capacity developed earlier.

\section{Whitney Stratification and Rank Geometry}
\label{app:whitney}

Many of the spaces arising in representation learning are not smooth manifolds but are arranged into strata determined by singularities in Jacobians or fibers of realization maps. A rigorous treatment requires recalling the structure of Whitney stratifications. Let $X$ be a semi-algebraic subset of $\mathbb{R}^n$. A Whitney stratification of $X$ is a decomposition into smooth submanifolds $\{S_\alpha\}$ that satisfy Whitney's conditions (A) and (B), the latter being crucial for ensuring that tangent spaces behave predictably near boundaries of strata. These conditions guarantee that analytic or semi-analytic functions defined on $X$ exhibit controlled behavior across transitions between strata.

In the context of parameterized models, the Jacobian $\ParamJac_\theta$ of the realization map has rank that varies across $\ThetaSpace$, and this variable rank induces the structure of strata. The stability theorems of Mather and Thom show that such stratifications persist under perturbations and that mappings between stratified spaces can preserve the stratified structure when certain transversality conditions hold. These results justify treating the movement across strata as a form of representational phase transition in agency.

\section{Operator-Theoretic Structure of Lamphrodynamic Energies}
\label{app:operator-theory}

Lamphrodynamic energies arise from differential operators acting on semantic fields. To formulate these energies precisely, it is necessary to view the underlying space of functions as a Hilbert space equipped with appropriate inner products induced by input distributions. Let $\mathcal{H} = L^2(\Omega,\mu)$ be the Hilbert space of real-valued functions on a domain $\Omega$ with respect to an input measure $\mu$. The gradient, Laplacian, and bi-Laplacian operators act densely on $\mathcal{H}$, and under mild conditions on $\Omega$, the operator $H = \alpha(-\Delta) + \beta(-\Delta)^2$ is self-adjoint and positive. This ensures the existence of a well-defined semigroup $e^{-tH}$ that governs lamphrodynamic smoothing.

The spectral decomposition of $H$ reveals how lamphrodynamic flows suppress high-frequency components of a semantic field. Eigenfunctions corresponding to large eigenvalues are attenuated rapidly, whereas those associated with small eigenvalues persist. This spectral behavior parallels classical results in heat flow and biharmonic diffusion and provides a principled basis for interpreting lamphrodynamic energy as a regularizer that promotes coherence.

The algebra generated by differential and multiplication operators forms a noncommutative algebra encoding the permissible internal semantic operations. Cohomological invariants of this algebra capture the semantic relations that remain stable under all transformations. When applied to RSVP, these invariants represent semantic attractors and long-lived conceptual structures.

\section{Variational Methods and Free-Energy Geometry}
\label{app:free-energy-geometry}

The variational free-energy functional plays a central role in both classical and generalized Bayesian inference. When viewed through the RSVP lens, free energy becomes a stratified functional whose geometry depends on the rank of the realization map. Let $(\theta,\phi)$ denote parameters of a generative model and a variational family. The free energy
\[
\FreeEnergy(\theta,\phi;x)
= \mathbb{E}_{q_\phi(z\mid x)}\!\left[\log q_\phi(z\mid x) - \log p_\theta(x,z)\right]
\]
is differentiable on each smooth stratum of the joint parameter space. The Hessian of $\FreeEnergy$ may be degenerate or ill-conditioned across strata, reflecting the collapse of semantic degrees of freedom. Morse theory adapted to stratified spaces, as developed by Goresky and MacPherson, allows for an analysis of critical points even in the presence of singularities.

The free-energy gradient flow interacts with the RSVP fields through induced metrics and curvature terms. Agents move toward configurations that minimize representational tension while remaining constrained by available resources. This geometric interpretation situates active inference within a broader class of lamphrodynamic systems.

\section{Renormalization of Priors and Multi-Scale Structure}
\label{app:renormalization}

Renormalization describes how semantic and probabilistic structures evolve across scales. Let $\Theta_\ell$ denote a representational manifold at resolution $\ell$, and let $p_\ell$ be the corresponding prior. Coarse-graining maps $\pi_\ell : \Theta_\ell \to \Theta_{\ell+1}$ induce renormalized priors through
\[
p_{\ell+1}(\theta_{\ell+1})
\propto
\int_{\pi_\ell^{-1}(\theta_{\ell+1})}
\exp\big(-\mathcal{C}_\ell(\theta)\big)\,d\mu_\ell(\theta),
\]
where $\mathcal{C}_\ell$ is the scale-dependent cost functional encoding curvature, entropy, and energy expenditure.

Under suitable conditions, the renormalized priors remain within the same exponential family, exhibiting Gibbs stability. The evolution of such priors defines a flow in the space of information-geometric metrics. Negative curvature accelerates the loss of fine-scale distinctions, while positive curvature preserves or even amplifies structure. The resulting semantic horizon bounds the depth of viable inference and establishes the limit at which distinctions remain operationally meaningful.

\section{Computational Mortality and Agency Collapse}
\label{app:computational-mortality}

Finally, computational mortality emerges as a consequence of finite energy budgets and dissipation. Let $E(t)$ denote the available computational energy. The dynamics of $E(t)$ obey
\[
\frac{dE}{dt} = -P(t) + I(t),
\]
where $P(t)$ captures computational effort and $I(t)$ represents replenishment. Mortality occurs when $E(t)$ reaches a minimum threshold below which the system cannot sustain coherent computation.

Within RSVP, computational mortality manifests as a collapse of strata: high-rank strata become inaccessible, the semantic manifold contracts, and agency reduces to reflexive or low-dimensional behavior. The study of optimal control policies that maximize the mortality time thus becomes integral to an analysis of sustainable intelligent behavior.

\newpage
\printbibliography

\end{document}
